% Options for packages loaded elsewhere
\PassOptionsToPackage{unicode}{hyperref}
\PassOptionsToPackage{hyphens}{url}
\PassOptionsToPackage{dvipsnames,svgnames,x11names}{xcolor}
%
\documentclass[
  letterpaper,
  DIV=11,
  numbers=noendperiod]{scrartcl}

\usepackage{amsmath,amssymb}
\usepackage{lmodern}
\usepackage{iftex}
\ifPDFTeX
  \usepackage[T1]{fontenc}
  \usepackage[utf8]{inputenc}
  \usepackage{textcomp} % provide euro and other symbols
\else % if luatex or xetex
  \usepackage{unicode-math}
  \defaultfontfeatures{Scale=MatchLowercase}
  \defaultfontfeatures[\rmfamily]{Ligatures=TeX,Scale=1}
  \setmainfont[]{Times New Roman}
  \setmonofont[]{Times New Roman}
\fi
% Use upquote if available, for straight quotes in verbatim environments
\IfFileExists{upquote.sty}{\usepackage{upquote}}{}
\IfFileExists{microtype.sty}{% use microtype if available
  \usepackage[]{microtype}
  \UseMicrotypeSet[protrusion]{basicmath} % disable protrusion for tt fonts
}{}
\makeatletter
\@ifundefined{KOMAClassName}{% if non-KOMA class
  \IfFileExists{parskip.sty}{%
    \usepackage{parskip}
  }{% else
    \setlength{\parindent}{0pt}
    \setlength{\parskip}{6pt plus 2pt minus 1pt}}
}{% if KOMA class
  \KOMAoptions{parskip=half}}
\makeatother
\usepackage{xcolor}
\usepackage[top=30mm,left=30mm]{geometry}
\setlength{\emergencystretch}{3em} % prevent overfull lines
\setcounter{secnumdepth}{2}
% Make \paragraph and \subparagraph free-standing
\ifx\paragraph\undefined\else
  \let\oldparagraph\paragraph
  \renewcommand{\paragraph}[1]{\oldparagraph{#1}\mbox{}}
\fi
\ifx\subparagraph\undefined\else
  \let\oldsubparagraph\subparagraph
  \renewcommand{\subparagraph}[1]{\oldsubparagraph{#1}\mbox{}}
\fi


\providecommand{\tightlist}{%
  \setlength{\itemsep}{0pt}\setlength{\parskip}{0pt}}\usepackage{longtable,booktabs,array}
\usepackage{calc} % for calculating minipage widths
% Correct order of tables after \paragraph or \subparagraph
\usepackage{etoolbox}
\makeatletter
\patchcmd\longtable{\par}{\if@noskipsec\mbox{}\fi\par}{}{}
\makeatother
% Allow footnotes in longtable head/foot
\IfFileExists{footnotehyper.sty}{\usepackage{footnotehyper}}{\usepackage{footnote}}
\makesavenoteenv{longtable}
\usepackage{graphicx}
\makeatletter
\def\maxwidth{\ifdim\Gin@nat@width>\linewidth\linewidth\else\Gin@nat@width\fi}
\def\maxheight{\ifdim\Gin@nat@height>\textheight\textheight\else\Gin@nat@height\fi}
\makeatother
% Scale images if necessary, so that they will not overflow the page
% margins by default, and it is still possible to overwrite the defaults
% using explicit options in \includegraphics[width, height, ...]{}
\setkeys{Gin}{width=\maxwidth,height=\maxheight,keepaspectratio}
% Set default figure placement to htbp
\makeatletter
\def\fps@figure{htbp}
\makeatother

\usepackage{booktabs}
\usepackage{longtable}
\usepackage{array}
\usepackage{multirow}
\usepackage{wrapfig}
\usepackage{float}
\usepackage{colortbl}
\usepackage{pdflscape}
\usepackage{tabu}
\usepackage{threeparttable}
\usepackage{threeparttablex}
\usepackage[normalem]{ulem}
\usepackage{makecell}
\usepackage{xcolor}
\usepackage{amsmath}
\usepackage{caption}
\KOMAoption{captions}{tableheading}
\makeatletter
\makeatother
\makeatletter
\makeatother
\makeatletter
\@ifpackageloaded{caption}{}{\usepackage{caption}}
\AtBeginDocument{%
\ifdefined\contentsname
  \renewcommand*\contentsname{Table of contents}
\else
  \newcommand\contentsname{Table of contents}
\fi
\ifdefined\listfigurename
  \renewcommand*\listfigurename{List of Figures}
\else
  \newcommand\listfigurename{List of Figures}
\fi
\ifdefined\listtablename
  \renewcommand*\listtablename{List of Tables}
\else
  \newcommand\listtablename{List of Tables}
\fi
\ifdefined\figurename
  \renewcommand*\figurename{Figure}
\else
  \newcommand\figurename{Figure}
\fi
\ifdefined\tablename
  \renewcommand*\tablename{Table}
\else
  \newcommand\tablename{Table}
\fi
}
\@ifpackageloaded{float}{}{\usepackage{float}}
\floatstyle{ruled}
\@ifundefined{c@chapter}{\newfloat{codelisting}{h}{lop}}{\newfloat{codelisting}{h}{lop}[chapter]}
\floatname{codelisting}{Listing}
\newcommand*\listoflistings{\listof{codelisting}{List of Listings}}
\makeatother
\makeatletter
\@ifpackageloaded{caption}{}{\usepackage{caption}}
\@ifpackageloaded{subcaption}{}{\usepackage{subcaption}}
\makeatother
\makeatletter
\makeatother
\ifLuaTeX
  \usepackage{selnolig}  % disable illegal ligatures
\fi
\IfFileExists{bookmark.sty}{\usepackage{bookmark}}{\usepackage{hyperref}}
\IfFileExists{xurl.sty}{\usepackage{xurl}}{} % add URL line breaks if available
\urlstyle{same} % disable monospaced font for URLs
\hypersetup{
  pdftitle={International Security},
  pdfauthor={Kadir Jun Ayhan},
  colorlinks=true,
  linkcolor={blue},
  filecolor={Maroon},
  citecolor={Blue},
  urlcolor={Blue},
  pdfcreator={LaTeX via pandoc}}

\title{International Security}
\author{Kadir Jun Ayhan}
\date{1/15/23}

\begin{document}
\maketitle
\renewcommand*\contentsname{Contents}
{
\hypersetup{linkcolor=}
\setcounter{tocdepth}{3}
\tableofcontents
}
\hypertarget{syllabus-2023-spring}{%
\section*{Syllabus (2023 Spring)}\label{syllabus-2023-spring}}

\begin{longtable}[]{@{}
  >{\raggedright\arraybackslash}p{(\columnwidth - 4\tabcolsep) * \real{0.3056}}
  >{\raggedright\arraybackslash}p{(\columnwidth - 4\tabcolsep) * \real{0.4306}}
  >{\raggedright\arraybackslash}p{(\columnwidth - 4\tabcolsep) * \real{0.2639}}@{}}
\toprule()
\endhead
{Course Title} & International Security & {Course No.:} IS518 \\
{Department/ Major} & GSIS & {Credit/ Hours:} 3 \\
{Class time/ Classroom} & Tuesday 15:30 -- 18:15 / IEB 902 & \\
{Instructor} & {Name:} Kadir Jun Ayhan & {Department:} GSIS \\
& {E-mail:} \href{mailto:ayhan@ewha.ac.kr}{\nolinkurl{ayhan@ewha.ac.kr}}
& {Phone:} 02-3277-4628 \\
{Office Hours/ Office Location} &
\href{https://calendly.com/kadirayhan/officehours}{Make an appointment
here} & \\
\bottomrule()
\end{longtable}

\hypertarget{course-overview}{%
\section{\texorpdfstring{{\textbf{Course
Overview}}}{Course Overview}}\label{course-overview}}

\hypertarget{course-description}{%
\subsection{\texorpdfstring{\textbf{Course
Description}}{Course Description}}\label{course-description}}

The main target of this course is graduate students who are interested
in international security. The course aims to provide students with a
theoretical understanding of various aspects of international security.
The course begins with an overview of international relations theories,
and levels of analysis in its study. Throughout the course, we will
continue to discuss different international relations theories'
approaches to security, without being limited to only ``mainstream''
theories. More specifically, we will cover different strands of realism,
liberalism, neo-institutionalism, constructivism, the English School,
the Copenhagen School, poststructuralism, and feminism, among others. We
will delve into major concepts related to international security
beginning with security itself, followed by anarchy, and order. This
will be followed by discussions on the role of international
institutions and security communities for international security. We
will then discuss securitization theory, ontological security, and the
role of reputation and status for security. Finally, we will learn about
security in the digital age. There are many active learning exercises
embedded into the course throughout the semester, including a simulation
and games, to help visualize abstract theoretical concepts to facilitate
students' meaningful learning.

\hypertarget{prerequisites}{%
\subsection{\texorpdfstring{\textbf{Prerequisites}}{Prerequisites}}\label{prerequisites}}

None.

\hypertarget{course-format}{%
\subsection{\texorpdfstring{\textbf{Course
Format}}{Course Format}}\label{course-format}}

\begin{longtable}[]{@{}
  >{\raggedright\arraybackslash}p{(\columnwidth - 4\tabcolsep) * \real{0.2639}}
  >{\raggedright\arraybackslash}p{(\columnwidth - 4\tabcolsep) * \real{0.3750}}
  >{\raggedright\arraybackslash}p{(\columnwidth - 4\tabcolsep) * \real{0.3611}}@{}}
\toprule()
\endhead
{Lecture} & {Discussion/ Presentation} & {Experiment/ Practicum} \\
20\% & 40\% & 40\% \\
\bottomrule()
\end{longtable}

This course is designed to facilitate student-centric learning. The
assignments for this course are not merely for grading but designed in a
way to help students practice and improve their critical thinking,
researching, presentation and academic writing skills. There are no
in-class exams for this course. This is because, I believe that you put
so much effort in writing papers in exams to be read only by one person.
On the other hand, writing papers give students an opportunity to
research and keep their writings as working papers which can be improved
in the future.

Active participation in discussion is strongly encouraged for this
course. In order to motivate students to read all required articles and
come to class prepared for discussion, response papers and discussion
participation constitute 30\% of students' grade for this course. In the
beginning of every lecture, you will be asked to share what you have
written in your response papers to stimulate discussion (based on what
you have already prepared). We will discuss your response papers in the
first part of the class. The second part of the class will begin with my
summary and clarification of some of the arguments and concepts in the
assigned readings. Then, we will have one or two student presentations
(depending on the number of students) which will apply the theoretical
concepts in the readings to a case study. The requirements of the
assignments are detailed below.

I expect students to attend classes, to arrive on time, and to come to
class prepared to engage in class discussion by reading at least the
required readings, taking note of the key arguments, and identifying
strengths, weaknesses, and contradictions in these arguments.

\hypertarget{course-objectives}{%
\subsection{\texorpdfstring{\textbf{Course
Objectives}}{Course Objectives}}\label{course-objectives}}

Learning objectives:

\begin{itemize}
\tightlist
\item
  Recognizing various international relations theories' approaches to
  international security at different levels of analysis.
\item
  Comparing and interpreting strengths and weaknesses of different
  approaches to international security, as well as inferring their
  implications.
\item
  Applying theoretical understanding of international security to case
  studies (presentation), and to other new settings (games); and
  analyzing and critiquing these cases.
\end{itemize}

\hypertarget{evaluation-system}{%
\subsection{\texorpdfstring{\textbf{Evaluation
System}}{Evaluation System}}\label{evaluation-system}}

\begin{longtable}[]{@{}
  >{\raggedright\arraybackslash}p{(\columnwidth - 8\tabcolsep) * \real{0.2000}}
  >{\raggedright\arraybackslash}p{(\columnwidth - 8\tabcolsep) * \real{0.2000}}
  >{\raggedright\arraybackslash}p{(\columnwidth - 8\tabcolsep) * \real{0.2000}}
  >{\raggedright\arraybackslash}p{(\columnwidth - 8\tabcolsep) * \real{0.2000}}
  >{\raggedright\arraybackslash}p{(\columnwidth - 8\tabcolsep) * \real{0.2000}}@{}}
\toprule()
\begin{minipage}[b]{\linewidth}\raggedright
{Mid-term Paper}
\end{minipage} & \begin{minipage}[b]{\linewidth}\raggedright
{Final Paper}
\end{minipage} & \begin{minipage}[b]{\linewidth}\raggedright
{Response Papers}
\end{minipage} & \begin{minipage}[b]{\linewidth}\raggedright
{Presentation}
\end{minipage} & \begin{minipage}[b]{\linewidth}\raggedright
{Participation and Discussion}
\end{minipage} \\
\midrule()
\endhead
15\% & 40\% & 10\% & 15\% & 20\% \\
\bottomrule()
\end{longtable}

Explanation of the evaluation system:

\hypertarget{response-papers}{%
\subsubsection{\texorpdfstring{\textbf{\emph{Response
Papers}}}{Response Papers}}\label{response-papers}}

Length: Less than half a page.

In order to prepare for class discussion, you must read all required
articles and submit weekly response papers beginning from Week 6. Each
student needs to submit at least four response papers (that is for four
different weeks). The presenter(s) should not submit a response paper.
The weekly response paper consists of two parts. The first part will be
a short summary of required articles, linking them together based on
common themes. Make sure to include MAIN puzzles, and arguments. In the
second part, you will criticize the articles by identifying problems in
them. Weekly response papers need to be submitted to Cyber Campus before
class time. You are encouraged to, but you don't have to, make comments
to your colleagues' response papers.

\hypertarget{presentations}{%
\subsubsection{\texorpdfstring{\textbf{\emph{Presentations}}}{Presentations}}\label{presentations}}

Length: Around 15 minutes

All students will make one presentation on a case study. The presenter
must read the assigned readings carefully when preparing for the
presentation. You are asked to present a case study that is related to
the required readings. You must make use of the theories/ frameworks/
concepts in the assigned readings for your analysis of the case study of
your choice. For example, if the required article looks at NATO as a
security community, you can apply a similar approach to analyze the case
of ASEAN as a security community. Make sure that you are not merely
summarizing another article. The point of the presentation is for you to
apply what you learned from the readings to a case study. You must show
your analysis and your value-added. Discussion will follow presentations
(based on discussion questions prepared by the presenter), and all
students must actively participate in this discussion.

\hypertarget{sec-mid-term}{%
\subsubsection{\texorpdfstring{\textbf{\emph{Diplomacy Game and Mid-term
Paper}}}{Diplomacy Game and Mid-term Paper}}\label{sec-mid-term}}

Due date: 25 Apr 23:59

We will play the \href{https://www.playdiplomacy.com/}{Diplomacy board
game which simulates the European balance of power at the time of World
War I}. You will write a short paper (900-1100 words) connecting what
happened in the game with different theoretical approaches to
international security. You can find the guidelines for the essay
\href{mid-term-instructions.qmd}{here}. You will submit your mid-term to
Cyber Campus (Turnitin).

Please note that there will not be a class on 25 Apr in lieu of mid-term
exam week.

\hypertarget{sec-final}{%
\subsubsection{\texorpdfstring{\textbf{\emph{Final
Paper}}}{Final Paper}}\label{sec-final}}

Proposal date: 30 May Class time

Due date: 20 Jun 23:59

Word count: 4000-6000 words

You can select any topic that is relevant to international security. You
must choose your topics by Week 6 and begin your research. You must
consult with me for your topic of choice. There are two ways you can
write your final paper. You can conduct your own research and present
your findings on your topic. I don't expect a full-fledged journal
article, but your research must have a good research question,
literature review, and research design. Those who choose this can have a
foundation upon which you can build to make it a publishable piece
later. Alternatively, you can do a critical literature review on a
specific topic. For the critical literature review, you will summarize
and evaluate what has been written on your chosen topic. You will
identify common themes and contrasting views in different papers.
Finally, you will critically discuss completeness of the literature and
identify gaps for future research in that field. Students will propose
their final papers on 30 May to get feedback from me and from their
colleagues. By the time of the proposal, you must have a clear research
question, research design, and literature review. You will have seven
minutes for your proposals, which will be followed by Q/A and feedback.

Here are guidelines that you can refer to when writing your final paper:

\begin{itemize}
\tightlist
\item
  Ackerman, E. (2018). ``Analyze This'': Writing in the Social Sciences.
  In G. Graff, C. Birkenstein, \& R. Durst (Eds.), ``They Say, I Say'':
  The Moves that Matter in Academic Writing: With Readings
  (pp.~187-205). New York: W. W. Norton \& Company.
\item
  Gustafsson, K., \& Hagström, L. (2018). What is the point? Teaching
  graduate students how to construct political science research puzzles.
  European Political Science, 17(4), 634-648. (see especially 643-645)
\end{itemize}

These short guidelines on writing may also be helpful:

\href{https://docs.google.com/spreadsheets/d/1TP7h17NHJFGCWGl67xF7rFJLcOpZg6pn5uSeZCVfLgA/edit\#gid=0}{Crafting
Good Research Questions};
\href{https://deanramser.files.wordpress.com/2017/01/1summaryvsanalysis_anthro.pdf}{Summary
\& Description vs.~Analysis \& Argument};
\href{https://www.nottingham.ac.uk/studentservices/documents/description-vs-analysis---learnhigher.pdf}{Critical
Thinking and Reflection};
\href{https://intranet.birmingham.ac.uk/as/libraryservices/library/asc/documents/public/pgtcriticalwriting.pdf}{A
short guide to critical writing}.

If you want to write a publishable article (by improving your final
paper later), I strongly encourage you to have this book:

\begin{itemize}
\tightlist
\item
  \href{https://www.amazon.com/Writing-Journal-Twelve-Weeks-Second-dp-022649991X/dp/022649991X/ref=dp_ob_title_bk}{Belcher,
  W. L. (2019). Writing Your Journal Article in Twelve Weeks, Second
  Edition: A Guide to Academic Publishing Success (2nd ed.). Chicago:
  University of Chicago Press}. (First edition is fine as well.)
\end{itemize}

\hypertarget{bonus-points}{%
\subsubsection{\texorpdfstring{\textbf{\emph{Bonus
Points}}}{Bonus Points}}\label{bonus-points}}

There will be bonus points for the winners of the Diplomacy game. This
bonus point scheme and rules will be explained during the semester.

\hypertarget{course-materials-and-additional-readings}{%
\section{\texorpdfstring{{\textbf{Course Materials and Additional
Readings}}}{Course Materials and Additional Readings}}\label{course-materials-and-additional-readings}}

\hypertarget{required-materials}{%
\subsection{Required Materials}\label{required-materials}}

No required textbooks. Refer to the detailed course schedule below.

\hypertarget{supplementary-materials}{%
\subsection{Supplementary Materials}\label{supplementary-materials}}

Refer to the detailed course schedule below.

\hypertarget{optional-additional-readings}{%
\subsection{Optional Additional
Readings}\label{optional-additional-readings}}

Refer to the detailed course schedule below.

\hypertarget{course-policies}{%
\section{\texorpdfstring{{\textbf{Course
Policies}}}{Course Policies}}\label{course-policies}}

\hypertarget{citation-and-plagiarism}{%
\subsection{Citation and Plagiarism}\label{citation-and-plagiarism}}

Students must use proper citation and avoid plagiarism. Any citation
style is fine as long as it is consistent throughout the paper.
Plagiarism will not be tolerated and severely punished. The papers will
be uploaded to Turnitin (via Cyber Campus).

\href{(https://www.nyu.edu/about/policies-guidelines-compliance/policies-and-guidelines/academic-integrity-for-students-at-nyu.html)}{``Plagiarism:
presenting others' work without adequate acknowledgement of its source,
as though it were one's own. Plagiarism is a form of fraud. We all stand
on the shoulders of others, and we must give credit to the creators of
the works that we incorporate into products that we call our own. Some
examples of plagiarism:}

\begin{itemize}
\item
  \href{(https://www.nyu.edu/about/policies-guidelines-compliance/policies-and-guidelines/academic-integrity-for-students-at-nyu.html)}{a
  sequence of words incorporated without quotation marks;}
\item
  \href{(https://www.nyu.edu/about/policies-guidelines-compliance/policies-and-guidelines/academic-integrity-for-students-at-nyu.html)}{an
  unacknowledged passage paraphrased from another's word;}
\item
  \href{(https://www.nyu.edu/about/policies-guidelines-compliance/policies-and-guidelines/academic-integrity-for-students-at-nyu.html)}{the
  use of ideas, sound recordings, computer data or images created by
  others as though it were one's own.''}
\end{itemize}

See
\href{https://owl.purdue.edu/owl/research_and_citation/resources.html}{also
this link}.

\hypertarget{late-submissions}{%
\subsection{Late Submissions}\label{late-submissions}}

For late submissions, you will get 80\% of your grading unless you have
a valid excuse. You can always contact me if you have a valid excuse to
ask for extension. I do not require students to submit official
documents (doctor report etc.). Your word is enough for me.

\newpage{}

\hypertarget{course-schedule}{%
\section{\texorpdfstring{{Course
Schedule}}{Course Schedule}}\label{course-schedule}}

\begin{longtable}{rll}
\toprule
Week & Date & Topics \& Class Materials, Assignments \\ 
\midrule
1 & 07 Mar & Introduction of the Course \\ 
2 & 14 Mar & No class - (Make-up class later) \\ 
3 & 21 Mar & Overview of International Relations Theories \\ 
4 & 28 Mar & Diplomacy Game 1 \\ 
5 & 04 Apr & Diplomacy Game 2 and Debriefing \\ 
6 & 11 Apr & What is Security? \\ 
7 & 18 Apr & Anarchy, Order, and Security \\ 
8 & 25 Apr & Mid-term Week \\ 
9 & 02 May & International Institutions and Security \\ 
10 & 09 May & Security Communities \\ 
11 & 16 May & Securitization Theory \\ 
12 & 23 May & Ontological Security \\ 
13 & 30 May & Final Paper Proposals \\ 
14 & 06 Jun & Reputation and Status for Security \\ 
15 & 13 Jun & Cybersecurity \\ 
16 & 20 Jun & Submission of Final Papers \\ 
17 & 27 Jun &  \\ 
\bottomrule
\end{longtable}

\newpage{}

\hypertarget{detailed-course-schedule-with-readings}{%
\section{\texorpdfstring{{Detailed Course Schedule with
Readings}}{Detailed Course Schedule with Readings}}\label{detailed-course-schedule-with-readings}}

\hypertarget{week-1-introduction-of-the-course}{%
\subsection*{Week 1: Introduction of the
Course}\label{week-1-introduction-of-the-course}}
\addcontentsline{toc}{subsection}{Week 1: Introduction of the Course}

Introduction of the course including course contents, assignments, and
expectations.

\hypertarget{week-2-no-class---make-up-class-later}{%
\subsection*{Week 2: No class - (Make-up class
later)}\label{week-2-no-class---make-up-class-later}}
\addcontentsline{toc}{subsection}{Week 2: No class - (Make-up class
later)}

\hypertarget{week-3-overview-of-international-relations-theories}{%
\subsection*{Week 3: Overview of International Relations
Theories}\label{week-3-overview-of-international-relations-theories}}
\addcontentsline{toc}{subsection}{Week 3: Overview of International
Relations Theories}

\textbf{Required Readings:}

\begin{itemize}
\tightlist
\item
  Finnemore, M., \& Sikkink, K. (2001). {TAKING} {STOCK}: {The}
  {Constructivist} {Research} {Program} in {International} {Relations}
  and {Comparative} {Politics}. \emph{Annual Review of Political
  Science}, \emph{4}(1), 391--416.
  \url{https://doi.org/10.1146/annurev.polisci.4.1.391}
\item
  Oneal, J. R., \& Russett, B. (1999). The {Kantian} {Peace}: {The}
  {Pacific} {Benefits} of {Democracy}, {Interdependence}, and
  {International} {Organizations}, 1885-1992. \emph{World Politics},
  \emph{52}(1), 1--37. \url{https://doi.org/10.1017/S0043887100020013}
\item
  Waltz, K. N. (1993). The emerging structure of international politics.
  \emph{International Security}, \emph{18}(2), 44--79.
  \url{https://doi.org/10.2307/2539097}
\end{itemize}

\textbf{Recommended Readings:}

\begin{itemize}
\tightlist
\item
  Buzan, B. (2014). An Introduction to the English School of
  International Relations: The Societal Approach. Cambridge: Polity
  Press. (12-39)
\item
  Hobden, S., \& Jones, R. W. (2014). Marxist theories of international
  relations. In J. Baylis, S. Smith, \& P. Owens (Eds.), The
  globalization of world politics: An introduction to international
  relations (6th Edition ed., 141-154). Oxford: Oxford University Press.
\item
  Keohane, R. O., \& Nye, J. S. (2012 {[}1997{]}). Power and
  Interdependence (4th ed.). Boston: Longman, Chapters 1-3.
\item
  Mearsheimer, J. J. (2013). Structural Realism. In T. Dunne, M. Kurki,
  \& S. Smith (Eds.), International Relations Theories: Discipline and
  Diversity (3rd Edition ed., 77-93). Oxford: Oxford University Press.
\item
  Mingst, Karen and Arreguin-Toft, Ivan M. (2017). Essentials of
  International Relations (7th Edition), W.W. Norton Company, Ch. 3
\item
  Nau, H. R. (2019). Perspectives on international relations: power,
  institutions, and ideas (6th ed.). Thousand Oaks, California: CQ
  Press, 52-54, 151-161.
\item
  Singer, J. D. (1961). The Level-of-Analysis Problem in International
  Relations. World Politics, 14(1), 77-92. Simmons (Eds.), Handbook of
  International Relations (170-194). London: SAGE Publications.
\item
  Sjoberg, Laura, \& Tickner, J. Ann (2013). Feminist Perspectives on
  International Relations. In W. Carlsnaes, T. Risse, \& B. - A.
  Slaughter, Anne-Marie (2011). International Relations, Principal
  Theories. In R. Wolfrum (Ed.), Max Planck Encyclopedia of Public
  International Law.
\item
  Wendt, A. (1992). Anarchy is what states make of it: the social
  construction of power politics. International Organization, 46(02),
  391-425.
\item
  Zehfuss, Maja (2013). Critical Theory, Poststructuralism, and
  Postcolonialism. In W. Carlsnaes, T. Risse, \& B. A. Simmons (Eds.),
  Handbook of International Relations (145-169). London: SAGE
  Publications.
\end{itemize}

\hypertarget{week-4-diplomacy-game-1}{%
\subsection*{Week 4: Diplomacy Game 1}\label{week-4-diplomacy-game-1}}
\addcontentsline{toc}{subsection}{Week 4: Diplomacy Game 1}

Here are the \href{diplomacy-game.qmd}{guidelines} for the game, and the
\href{mid-term-instructions.qmd}{instructions for the essay} based on
the game.

\textbf{Recommended Readings:}

\begin{itemize}
\tightlist
\item
  Bull, H. (2012 {[}1977{]}). The Anarchical Society: A Study of Order
  in World Politics (4th ed.). London: Palgrave Macmillan.
\item
  Wendt, Alexander (1999). Social Theory of International Politics.
  Cambridge, UK: Cambridge University Press. Chapter 6.
\item
  Wight, Martin (1991). The Three Traditions of International Theory. In
  International Theory: The Three Traditions, eds.~G. Wight and N.
  Porter. Leicester: Leicester University Press, 7--24.
\end{itemize}

\hypertarget{week-5-diplomacy-game-2-and-debriefing}{%
\subsection*{Week 5: Diplomacy Game 2 and
Debriefing}\label{week-5-diplomacy-game-2-and-debriefing}}
\addcontentsline{toc}{subsection}{Week 5: Diplomacy Game 2 and
Debriefing}

In the latter half of the class, we will discuss what we have learned
from this active learning exercise in an interactive session.

\textbf{Recommended Readings:}

\begin{itemize}
\tightlist
\item
  Charles L. Glaser (1997). The Security Dilemma Revisited. World
  Politics, 50(1), 171-201.
\item
  Christensen, T. J., \& Snyder, J. (1990). Chain Gangs and Passed
  Bucks: Predicting Alliance Patterns in Multipolarity. International
  Organization, 44(2), 137-168. Retrieved from
  http://www.jstor.org/stable/2706792
\item
  Jervis, Robert (1978). Cooperation under the Security Dilemma. World
  Politics, 30(2), 167-214.
\item
  Leeds, Brett Ashley (2003). Alliance reliability in times of war:
  Explaining state decisions to violate treaties. International
  Organization, 57(4), 801-827.
\item
  Snyder, G. H. (1984). The Security Dilemma in Alliance Politics. World
  Politics, 36(4), 461-495.
\item
  Walt, S. M. (1985). Alliance Formation and the Balance of World Power.
  International Security, 9(4), 3-43.
\end{itemize}

\hypertarget{week-6-what-is-security}{%
\subsection*{Week 6: What is Security?}\label{week-6-what-is-security}}
\addcontentsline{toc}{subsection}{Week 6: What is Security?}

\textbf{Required Readings:}

\begin{itemize}
\tightlist
\item
  Burke, A. (2017). Security. In J. George, R. Devetak, \& S. Percy
  (Eds.), \emph{An {Introduction} to {International} {Relations}} (3rd
  ed., pp. 198--209). Cambridge University Press.
  \url{https://doi.org/10.1017/9781316855188}
\item
  Buzan, B., \& Hansen, L. (2009). \emph{The {Evolution} of
  {International} {Security} {Studies}}. Cambridge University Press.
  \url{https://doi.org/10.1017/CBO9780511817762}, Chapter 2.
\item
  Rothschild, E. (1995). What {Is} {Security}? \emph{Daedalus},
  \emph{124}(3), 53--98. \url{https://www.jstor.org/stable/20027310}
\end{itemize}

\textbf{Recommended Readings:}

\begin{itemize}
\tightlist
\item
  Barkawi, T., \& Laffey, M. (2006). The Postcolonial Moment in Security
  Studies. Review of International Studies, 32(2), 329-352.
\item
  Barkawi, Tarak (2016). Decolonising war. European Journal of
  International Security, 1(2), 199-214.
\item
  Browning, C. S. (2013). International security: A very short
  introduction. Oxford: Oxford University Press.
\item
  Buzan, Barry (2015). The English School: A neglected approach to
  international security studies. Security Dialogue, 46(2), 126-143.
\item
  Buzan, Barry, and Lene Hansen (2009). The Evolution of International
  Security Studies. Cambridge University Press. Chapters 1, 3.
\item
  Haftendorn, Helga (1991). The Security Puzzle: Theory-Building and
  Discipline-Building in International Security. International Studies
  Quarterly, 35(1), 3-17.
\item
  Hudson, Valerie M., Mary Caprioli, Bonnie Ballif-Spanvill, Rose
  McDermott, and Chad F. Emmett (2009). The heart of the matter: The
  security of women and the security of states. International security,
  33(3), 7-45.
\item
  Lebow, R. Ned (2010). Why Nations Fight: Past and Future Motives for
  War. Cambridge: Cambridge University Press. Chapter 1.
\item
  Levy, Jack S. (1998). The Causes of War and the Conditions of Peace.
  Annual Review of Political Science, 1(1), 139-165.
\item
  Paris, R. (2001). Human Security: Paradigm Shift or Hot Air?
  International Security, 26(2), 87-102.
\item
  Pasha, M. K. (1996). Security as Hegemony. Alternatives: Global,
  Local, Political, 21(3), 283-302.
\item
  Pratt, Nicola (2014). Reconceptualizing Gender, Reinscribing
  Racial--Sexual Boundaries in International Security: The Case of UN
  Security Council Resolution 1325 on ``Women, Peace and Security''.
  International Studies Quarterly, 57(4), 772-783.
\item
  Suganami, H. (2017). The causes of war. In J. George, R. Devetak, \&
  S. Percy (Eds.), An Introduction to International Relations (3rd ed.,
  pp.~224-234). Cambridge: Cambridge University Press.
\end{itemize}

\hypertarget{week-7-anarchy-order-and-security}{%
\subsection*{Week 7: Anarchy, Order, and
Security}\label{week-7-anarchy-order-and-security}}
\addcontentsline{toc}{subsection}{Week 7: Anarchy, Order, and Security}

\textbf{Required Readings:}

\begin{itemize}
\tightlist
\item
  Adler-Nissen, R., \& Zarakol, A. (2020). Struggles for recognition:
  {The} {Liberal} international order and the merger of its discontents.
  \emph{International Organization}, 1--24.
  \url{https://doi.org/10.1017/S0020818320000454}
\item
  Bull, H. (2012). \emph{The anarchical society: A study of order in
  world politics}. Bloomsbury Publishing.
  \url{https://doi.org/10.1007/978-1-349-24028-9}, Chapter 2.
\item
  Lebow, R. N. (2005). Power, {Persuasion} and {Justice}.
  \emph{Millennium}, \emph{33}(3), 551--581.
  \url{https://doi.org/10.1177/03058298050330031501}
\end{itemize}

\textbf{Recommended Readings:}

Recommended Readings:

\begin{itemize}
\tightlist
\item
  Acharya, Amitav (2017). After liberal hegemony: The advent of a
  multiplex world order. Ethics \& International Affairs, 31(3),
  271-285.
\item
  Barnett, Michael (2020). COVID-19 and the Sacrificial International
  Order. International Organization, 1-20.
\item
  Brooks, S. G., \& Wohlforth, W. C. (2016). The Rise and Fall of the
  Great Powers in the Twenty-first Century: China's Rise and the Fate of
  America's Global Position. International security, 40(3), 7-53.
  https://doi.org/10.1162/ISEC\_a\_00225
\item
  Finnemore, Martha (2009). Legitimacy, hypocrisy, and the social
  structure of unipolarity. World Politics, 61(1), 58-85.
\item
  Goddard, Stacie E., \& Nexon, Daniel H. (2016). The Dynamics of Global
  Power Politics: A Framework for Analysis. Journal of Global Security
  Studies, 1(1).
\item
  Hall, John A. (1996). International orders. Cambridge, MA: Polity
  Press. Chapter 1.
\item
  Hobson, J. M. (2014). The twin self-delusions of IR: why `hierarchy'
  and not `anarchy' is the core concept of IR. Millennium, 42(3),
  557-575. -Ikenberry, John (2001). After Victory: Institutions,
  Strategic Restraint, and the Rebuilding of Order After Major Wars.
  Princeton, NJ: Princeton University Press
\item
  Lebow, R. Ned (2008). A Cultural Theory of International Relations.
  Cambridge: Cambridge University Press. Chapters 1, 2, 3, 10.
\item
  Park, Seo-Hyun (2017). Sovereignty and Status in East Asian
  International Relations. Cambridge: Cambridge University Press.
\item
  Wohlforth, W. C. (1999). The Stability of a Unipolar World.
  International Security, 24(1), 5-41.
\end{itemize}

\hypertarget{week-8-mid-term-week}{%
\subsection*{Week 8: Mid-term Week}\label{week-8-mid-term-week}}
\addcontentsline{toc}{subsection}{Week 8: Mid-term Week}

See Mid-term Section (Section~\ref{sec-mid-term}) for details.

\hypertarget{week-9-international-institutions-and-security}{%
\subsection*{Week 9: International Institutions and
Security}\label{week-9-international-institutions-and-security}}
\addcontentsline{toc}{subsection}{Week 9: International Institutions and
Security}

\textbf{Required Readings:}

\begin{itemize}
\tightlist
\item
  Cui, S., \& Buzan, B. (2016). Great {Power} {Management} in
  {International} {Society}. \emph{The Chinese Journal of International
  Politics}, \emph{9}(2), 181--210.
  \url{https://doi.org/10.1093/cjip/pow005}
\item
  Ikenberry, G. J. (2001). \emph{After {Victory}: {Institutions},
  {Strategic} {Restraint}, and the {Rebuilding} of {Order} {After}
  {Major} {Wars}}. Princeton University Press.
  \url{https://doi.org/10.2307/j.ctv3znx0v}, Chapter 2.
\end{itemize}

\textbf{Recommended Readings:}

\begin{itemize}
\tightlist
\item
  Boehmer, C., Gartzke, E., \& Nordstrom, T. (2004). Do
  intergovernmental organizations promote peace? World Politics, 57(1),
  1-38.
\item
  Jervis, R. (1982). Security Regimes. International Organization,
  36(2), 357-378.
\item
  Lake, David A. (2001). Beyond anarchy - The importance of security
  institutions. International Security, 26(1), 129-160.
\item
  Mearsheimer, J. J. (1995). The False Promise of International
  Institutions. International Security, 19(3), 5-49.
\item
  Pouliot, V. (2016). Hierarchy in Practice: Multilateral Diplomacy and
  the Governance of International Security, European Journal of
  International Security, 1(1), 5-26.
\item
  Ruggie, J. G. (1982). International regimes, transactions, and change:
  embedded liberalism in the postwar economic order. International
  Organization, 36(2), 379-415.
\end{itemize}

\hypertarget{week-10-security-communities}{%
\subsection*{Week 10: Security
Communities}\label{week-10-security-communities}}
\addcontentsline{toc}{subsection}{Week 10: Security Communities}

\textbf{Required Readings:}

\begin{itemize}
\tightlist
\item
  Adler, E., Barnett, M., \& Smith, S. (1998). \emph{Security
  communities}. Cambridge University Press.
  \url{https://doi.org/10.1017/CBO9780511598661}
\item
  Crawford, N. C. (1994). A security regime among democracies:
  Cooperation among {Iroquois} nations. \emph{International
  Organization}, \emph{48}(3), 345--385.
  \url{https://doi.org/10.1017/S002081830002823X}
\item
  Hunter, R. (2022). The {Ukraine} {Crisis}: Why and what now?
  \emph{Survival}, \emph{64}(1), 7--28.
  \url{https://doi.org/10.1080/00396338.2022.2032953}
\end{itemize}

\textbf{Recommended Readings:}

\begin{itemize}
\tightlist
\item
  Adler, Emanuel, \& Greve, Patricia (2009). When Security Community
  Meets Balance of Power: Overlapping Regional Mechanisms of Security
  Governance. Review of International Studies, 35, 59-84.
\item
  Adler, Emanuel and Michael N. Barnett (1996). Governing Anarchy: A
  Research Agenda for the Study of Security Communities. Ethics \&
  International Affairs, 10, 63-98.
\item
  Bially Mattern, Janice (2005). Ordering International Politics
  Identity, Crisis, and Representational Force. New York: Routledge.
  Chapter 1.
\item
  Bially Mattern, Janice (2005). Why `Soft Power' Isn't So Soft:
  Representational Force and the Sociolinguistic Construction of
  Attraction in World Politics. Millennium - Journal of International
  Studies, 33(3), 583-612.
\item
  Buzan, B., \& Waever, O. (2003). Regions and powers: the structure of
  international security. Cambridge: Cambridge University Press.
  Chapters 1-3.
\item
  Koschut, S. (2014). Emotional (security) communities: the significance
  of emotion norms in inter-allied conflict management. Review of
  International Studies, 40(3), 533-558.
\item
  Pouliot, Vincent (2008). The Logic of Practicality: A Theory of
  Practice of Security Communities. International Organization, 62(2),
  257-288.
\end{itemize}

\hypertarget{week-11-securitization-theory}{%
\subsection*{Week 11: Securitization
Theory}\label{week-11-securitization-theory}}
\addcontentsline{toc}{subsection}{Week 11: Securitization Theory}

\textbf{Required Readings:}

\begin{itemize}
\tightlist
\item
  Goddard, S. E., \& Krebs, R. R. (2015). Rhetoric, {Legitimation}, and
  {Grand} {Strategy}. \emph{Security Studies}, \emph{24}(1), 5--36.
  \url{https://doi.org/10.1080/09636412.2014.1001198}
\item
  Kirk, J. (2020). From {Threat} to {Risk}? {Exceptionalism} and
  {Logics} of {Health} {Security}. \emph{International Studies
  Quarterly}, \emph{64}(2), 266--276.
  \url{https://doi.org/10.1093/isq/sqaa021}
\item
  McDonald, M. (2008). Securitization and the {Construction} of
  {Security}. \emph{European Journal of International Relations},
  \emph{14}(4), 563--587. \url{https://doi.org/10.1177/1354066108097553}
\end{itemize}

\textbf{Recommended Readings:}

\begin{itemize}
\tightlist
\item
  Balzacq, T., Guzzini, S., Williams, M. C., Wæver, O., \& Patomäki, H.
  (2014). What kind of theory -- if any -- is securitization?
  International Relations, 29(1), 96-96.
\item
  Balzacq, T. (2015). The `Essence' of securitization: Theory, ideal
  type, and a sociological science of security. International Relations,
  29(1), 103-113.
\item
  Buzan, B., Wæver, O., \& de Wilde, J. (1998). Security: A New
  Framework for Analysis. Lynne Rienner. Chapters 1-2.
\item
  Lipscy, P. Y. (2020). COVID-19 and the Politics of Crisis.
  International Organization, 1-30.
\item
  Vaughan-Williams, N., \& Stevens, D. (2015). Vernacular theories of
  everyday (in)security: The disruptive potential of non-elite
  knowledge. Security Dialogue, 47(1), 40-58.
\item
  Wæver, O. (2011). Politics, security, theory. Security Dialogue,
  42(4-5), 465-480.
\end{itemize}

\hypertarget{week-12-ontological-security}{%
\subsection*{Week 12: Ontological
Security}\label{week-12-ontological-security}}
\addcontentsline{toc}{subsection}{Week 12: Ontological Security}

\textbf{Required Readings:}

\begin{itemize}
\tightlist
\item
  Mitzen, J. (2006). Ontological {Security} in {World} {Politics}:
  {State} {Identity} and the {Security} {Dilemma}. \emph{European
  Journal of International Relations}, \emph{12}(3), 341--370.
  \url{https://doi.org/10.1177/1354066106067346}
\item
  Subotić, J. (2016). Narrative, {Ontological} {Security}, and {Foreign}
  {Policy} {Change} 1. \emph{Foreign Policy Analysis}, \emph{12}(4),
  610--627. \url{https://doi.org/10.1111/fpa.12089}
\item
  Zarakol, A. (2010). Ontological (in) security and state denial of
  historical crimes: {Turkey} and {Japan}. \emph{International
  Relations}, \emph{24}(1), 3--23.
  \url{https://doi.org/10.1177/0047117809359040}
\end{itemize}

\textbf{Recommended Readings:}

\begin{itemize}
\tightlist
\item
  Browning, Christopher S., \& Joenniemi, Pertti (2016). Ontological
  security, self-articulation and the securitization of identity.
  Cooperation and Conflict, 52(1), 31-47.
\item
  Kinnvall, Catarina (2016). Feeling ontologically (in)secure: States,
  traumas and the governing of gendered space. Cooperation and Conflict,
  52(1), 90-108.
\item
  Lebow, R. N. (2016). National identities and international relations.
  Cambridge: Cambridge University Press. Chapter 2.
\item
  Pacher, A. (2019). The diplomacy of post-Soviet de facto states:
  ontological security under stigma. International Relations, 33(4),
  563-585.
\item
  Pratt, S. F. (2017). A relational view of ontological security in
  international relations. International Studies Quarterly, 61(1),
  78-85.
\item
  Zarakol, Ayşe (2016). States and ontological security: A historical
  rethinking. Cooperation and Conflict, 52(1), 48-68.
\end{itemize}

\hypertarget{week-13-final-paper-proposals}{%
\subsection*{Week 13: Final Paper
Proposals}\label{week-13-final-paper-proposals}}
\addcontentsline{toc}{subsection}{Week 13: Final Paper Proposals}

See Final Paper Section (Section~\ref{sec-final}) for details.

\hypertarget{week-14-reputation-and-status-for-security}{%
\subsection*{Week 14: Reputation and Status for
Security}\label{week-14-reputation-and-status-for-security}}
\addcontentsline{toc}{subsection}{Week 14: Reputation and Status for
Security}

\textbf{Required Readings:}

\begin{itemize}
\tightlist
\item
  Larson, D. W., \& Shevchenko, A. (2010). Status seekers: {Chinese} and
  {Russian} responses to {US} primacy. \emph{International Security},
  \emph{34}(4), 63--95. \url{https://doi.org/10.1162/isec.2010.34.4.63}
\item
  Lupton, D. L. (2018). Reexamining {Reputation} for {Resolve}:
  {Leaders}, {States}, and the {Onset} of {International} {Crises}.
  \emph{Journal of Global Security Studies}, \emph{3}(2), 198--216.
  \url{https://doi.org/10.1093/jogss/ogy004}
\item
  Renshon, J. (2016). Status {Deficits} and {War}. \emph{International
  Organization}, \emph{70}(3), 513--550.
  \url{https://doi.org/10.1017/S0020818316000163}
\end{itemize}

\textbf{Recommended Readings:}

\begin{itemize}
\tightlist
\item
  Dafoe, A., Renshon, J., \& Huth, P. (2014). Reputation and Status as
  Motives for War. Annual Review of Political Science, 17(1), 371-393.
\item
  de Carvalho, B., \& Neumann, I. B. (2014). Small State Status Seeking:
  Norway's Quest for International Standing: Taylor \& Francis.
\item
  Duque, M. G. (2018). Recognizing International Status: A Relational
  Approach. International Studies Quarterly, 62(3), 577-592.
  https://doi.org/10.1093/isq/sqy001
\item
  Jervis, Robert, Yarhi-Milo, Keren, \& Casler, Don (2020). Redefining
  the Debate Over Reputation and Credibility in International Security:
  Promises and Limits of New Scholarship. World Politics, 1-37.
\item
  Mercer, J. (2013). Emotion and Strategy in the Korean War,
  International Organization, 67(2), 221-252.
\item
  Park, Seo-Hyun (2017). Sovereignty and Status in East Asian
  International Relations. Cambridge: Cambridge University Press.
\item
  Paul, T. V., Larson, D. W., \& Wohlforth, W. C. (2014). Status in
  world politics. Cambridge: Cambridge University Press.
\item
  Wohlforth, W. C. (2009). Unipolarity, status competition, and great
  power war. World Politics, 61(1), 28-57.
\item
  Yarhi-Milo, Keren (2018). Who fights for reputation: the psychology of
  leaders in international conflict. Princeton University Press.
  Introduction chapter.
\end{itemize}

\hypertarget{week-15-cybersecurity}{%
\subsection*{Week 15: Cybersecurity}\label{week-15-cybersecurity}}
\addcontentsline{toc}{subsection}{Week 15: Cybersecurity}

There will be an invited speaker for this class. I will give the
required readings list later in the semester.

\textbf{Recommended Readings:}

\begin{itemize}
\tightlist
\item
  Branch, Jordan (2021). What's in a Name? Metaphors and Cybersecurity.
  International Organization, 75(1), 39-70.
\item
  Choucri, N., \& Goldsmith, D. (2012). Lost in cyberspace: Harnessing
  the Internet, international relations, and global security. Bulletin
  of the Atomic Scientists, 68(2), 70-77.
\item
  Crilley, Rhys, \& Chatterje-Doody, Precious. (2019). Security studies
  in the age of `post-truth' politics: in defence of poststructuralism.
  Critical Studies on Security, 7(2), 166-170.
  https://doi.org/10.1080/21624887.2018.1441634
\item
  Jackson, Susan T., Crilley, Rhys, Manor, Ilan, Baker, Catherine,
  Oshikoya, Modupe, Joachim, Jutte, Robinson, Nick, Schneiker, Andrea,
  Grove, Nicole Sunday, \& Enloe, Cynthia (2021). Militarization 2.0:
  communication and the normalization of political violence in the
  digital age. International Studies Review, 23(3), 1046-1071.
\item
  La Cour, Christina (2020). Theorising digital disinformation in
  international relations. International Politics, 57(4), 704-723.
\item
  Lanoszka, A. (2019). Disinformation in international politics.
  European Journal of International Security, 4(2), 227-248.
\item
  Nye Jr, J. S. (2016). Deterrence and dissuasion in cyberspace.
  International security, 41(3), 44-71.
\item
  Reinsberg, B. (2021). Fully-automated liberalism? Blockchain
  technology and international cooperation in an anarchic world.
  International Theory, 13(2), 287-313.
\end{itemize}

\hypertarget{week-16-submission-of-final-papers}{%
\subsection*{Week 16: Submission of Final
Papers}\label{week-16-submission-of-final-papers}}
\addcontentsline{toc}{subsection}{Week 16: Submission of Final Papers}

See Final Paper Section (Section~\ref{sec-final}) for details.

\hypertarget{special-accommodations}{%
\section{\texorpdfstring{{Special
Accommodations}}{Special Accommodations}}\label{special-accommodations}}

\begin{itemize}
\tightlist
\item
  According to the University regulation section \#57-3, students with
  disabilities can request for special accommodations related to
  attendance, lectures, assignments, or tests by contacting the course
  professor at the beginning of semester. Based on the nature of the
  students' request, students can receive support for such
  accommodations from the course professor or from the Support Center
  for Students with Disabilities (SCSD). Please refer to the below
  examples of the types of support available in the lectures,
  assignments, and evaluations.
\end{itemize}

\begin{longtable}[]{@{}
  >{\raggedright\arraybackslash}p{(\columnwidth - 4\tabcolsep) * \real{0.2917}}
  >{\raggedright\arraybackslash}p{(\columnwidth - 4\tabcolsep) * \real{0.2639}}
  >{\raggedright\arraybackslash}p{(\columnwidth - 4\tabcolsep) * \real{0.4444}}@{}}
\toprule()
\begin{minipage}[b]{\linewidth}\raggedright
Lecture
\end{minipage} & \begin{minipage}[b]{\linewidth}\raggedright
Assignments
\end{minipage} & \begin{minipage}[b]{\linewidth}\raggedright
Evaluation
\end{minipage} \\
\midrule()
\endhead
Visual impairment: braille, enlarged reading materials. Hearing
impairment: note-taking assistant. Physical impairment: access to
classroom, note-taking assistant. & Extra days for submission,
alternative assignments. & Visual impairment: braille examination paper,
examination with voice support, longer examination hours, note-taking
assistant. Hearing impairment: written examination instead of oral.
Physical impairment: longer examination hours, note-taking assistant. \\
\bottomrule()
\end{longtable}

\begin{itemize}
\item
  Actual support may vary depending on the course.
\item
  If you have other special needs, please let me know. I will do my best
  to flexibly accommodate your needs.
\end{itemize}

\hypertarget{notes}{%
\section*{\texorpdfstring{{Notes}}{Notes}}\label{notes}}
\addcontentsline{toc}{section}{{Notes}}

The contents of this syllabus are not final. I may update them later.

\hypertarget{read-the-syllabus}{%
\section*{\texorpdfstring{{Read the
Syllabus!}}{Read the Syllabus!}}\label{read-the-syllabus}}
\addcontentsline{toc}{section}{{Read the Syllabus!}}

\url{https://www.youtube.com/watch?v=xQJ1vQ4RJ5I\&t=1s}

\includegraphics{http://www.phdcomics.com/comics/archive/phd051013s.gif}



\end{document}
